\section[Capítulo 3: Gamificación y Juegos Serios]{\centering Capítulo 3: Gamificación y Juegos Serios}

La gamificación y los juegos serios permiten estructurar una experiencia educativa con reglas, retos, progreso y retroalimentación. En este proyecto, estos elementos se usan para ordenar la práctica de ciberseguridad y no solamente para hacer que la actividad parezca entretenida (Deterding et al., 2011).

Este capítulo desarrolla la relación entre juego y aprendizaje. La idea principal es que una mecánica de juego tiene valor educativo cuando obliga al estudiante a aplicar un criterio relacionado con el objetivo de formación (Kapp, 2012).

\subsection{Gamificación}

La gamificación consiste en utilizar elementos de diseño de juego en contextos que no son principalmente lúdicos. Entre estos elementos se encuentran niveles, misiones, puntos, progreso, retos, restricciones de tiempo y recompensas (Deterding et al., 2011).

En el videojuego propuesto, la gamificación aparece en el mapa de casos, el turno del analista, la reputación, los módulos activos y el reporte de cierre. Estos elementos ayudan a organizar la dificultad y a mantener la atención del estudiante sobre las decisiones que toma (Kapp, 2012).

\subsection{Juegos serios}

Los juegos serios son aplicaciones interactivas diseñadas con una finalidad formativa, de entrenamiento o concientización. Su objetivo no es eliminar la diversión, sino hacer que las mecánicas del juego sirvan a un propósito educativo definido (Michael \& Chen, 2006).

\begin{longtable}{|p{0.30\textwidth}|p{0.55\textwidth}|}
\hline
\textbf{Elemento de juego} & \textbf{Función educativa} \\
\hline
Rol de analista SOC & Sitúa al estudiante en un contexto profesional simulado (Michael \& Chen, 2006). \\
\hline
Casos progresivos & Organiza amenazas por dificultad y tema (Kapp, 2012). \\
\hline
Decisiones & Obliga a aplicar criterio, no solo recordar conceptos (Deterding et al., 2011). \\
\hline
Reporte final & Presenta consecuencias y permite revisar errores (Hattie \& Timperley, 2007). \\
\hline
\end{longtable}

En ciberseguridad, los juegos serios son pertinentes porque permiten representar ataques e incidentes sin exponer sistemas reales. Esta característica es útil para entrenar identificación de amenazas, priorización y respuesta dentro de un ambiente controlado (Ng \& Hasan, 2024).

\subsection{Relación con la enseñanza de ciberseguridad}

La relación entre gamificación y ciberseguridad se sostiene en que muchos incidentes pueden convertirse en escenarios de decisión. El estudiante recibe información, detecta señales, compara datos y elige una acción, lo cual se aproxima al razonamiento usado en análisis de seguridad (Irvine, 2005).

Antecedentes como CyberCIEGE muestran que los juegos pueden representar decisiones de seguridad informática dentro de entornos simulados. En el presente proyecto, esta idea se adapta a una experiencia más accesible: correos sospechosos, consola, llamadas, registro de turno y reportes de consecuencias (Irvine, 2005).

\subsection{Síntesis del capítulo}

La gamificación aporta estructura y motivación, mientras que los juegos serios aportan finalidad educativa. Juntas, ambas bases sostienen la propuesta de un videojuego que permita practicar ciberseguridad mediante observación, toma de decisiones y retroalimentación integrada al escenario (Deterding et al., 2011).

\newpage
